\chapter*{Graduate-Level Scope}
\addcontentsline{toc}{chapter}{Graduate-Level Scope}

\section*{Prerequisites}

The reader should be able to work with sets, functions, relations, finite
strings, countable sets, induction on natural numbers, and ordinary rigorous
proofs.  Prior exposure to elementary logic is useful but not
required.  Chapters that use cardinal arithmetic, algebra, or topology state
the additional facts they require.

\section*{Learning outcomes}

A reader completing the reference should be able to:
\begin{enumerate}
  \item define a formal language by an inductive grammar and prove results by
        structural induction;
  \item define capture-avoiding substitution and prove a substitution lemma;
  \item prove soundness for a natural-deduction or sequent calculus;
  \item prove completeness for classical propositional logic and explain a
        Henkin proof of completeness for classical first-order logic;
  \item derive compactness and the downward and upward L\"owenheim--Skolem
        theorems with their cardinal hypotheses;
  \item use back-and-forth arguments, elementary chains, and ultraproducts;
  \item prove undecidability by a many-one reduction;
  \item state and apply the diagonal lemma and derive the first incompleteness
        theorem, the second incompleteness theorem, and undefinability of truth
        with exact hypotheses;
  \item explain cut elimination, normalization, and the subformula property;
  \item give relational semantics for intuitionistic and normal modal logic and
        prove representative soundness results;
  \item distinguish full semantics from Henkin semantics for second-order logic;
        and
  \item audit a written or mechanized proof for hidden assumptions, variable
        capture, rule misuse, and shifts between object language and
        metalanguage.
\end{enumerate}

\tableofcontents

\mainmatter

